% ===================== Chapter 23 =====================
\chapter{Advanced Topics and Best Practices}
\label{ch:advanced}

You now know the whole language: types, control flow, functions, data structures,
generics, interfaces, closures, the standard library, concurrency, and the FFI.
This final chapter steps back from features to \emph{judgement} --- how to use what
you know well. It gathers advice on performance, structure, style, and the
pitfalls worth knowing, so that the programs you write next are not just correct
but good.

\section{Performance: fast by default, faster with care}

Salam programs are fast because they compile to C and run as native code, with
static types and no hidden boxing. For most programs that is the end of the
performance story --- you get speed for free. When you do need to squeeze more, a
few principles matter far more than micro-tricks:

\begin{itemize}[leftmargin=1.4em]
  \item \textbf{Choose the right data structure.} An algorithm that looks up by
        key thousands of times wants a \code{HashMap}, not a linear scan through a
        \code{Vector}. The biggest speed wins come from structure, not from
        shaving instructions.
  \item \textbf{Avoid needless allocation.} Building a string by concatenating in
        a tight loop allocates on every step. When you can, compute sizes up front
        or reuse a buffer. Reserve a vector's capacity if you know roughly how many
        elements it will hold.
  \item \textbf{Measure before optimising.} Intuition about what is slow is
        usually wrong. Time the program, find the part that actually dominates,
        and improve that --- not the part that merely \emph{looks} expensive.
  \item \textbf{Prefer the compiled path.} \code{salam exec} (the interpreter) is
        perfect for learning and quick runs, but a \code{salam build} executable
        is far faster for real work.
\end{itemize}

\begin{deep}[In Depth: optimization levels and the LLVM backend]
Because the default backend emits C, the C compiler's own optimizer does a great
deal of work on your behalf --- inlining small functions, unrolling loops,
eliminating dead code. Salam also has an experimental LLVM backend
(\code{salam llvm}) with explicit optimization levels (\code{-O0} through
\code{-O3}, plus \code{-Os}/\code{-Oz} for size), which can target native code or
WebAssembly. For everyday work the defaults are excellent; reach for the
optimization flags only when a measurement tells you to.
\end{deep}

\section{Structuring a program}

Small programs fit in one file; real ones do not. As a program grows, structure is
what keeps it understandable:

\begin{itemize}[leftmargin=1.4em]
  \item \textbf{Split by responsibility.} Put related types and functions together
        in a module (a file with a \code{package}), and keep each module focused on
        one area --- parsing, storage, the user interface. A reader should be able
        to guess what is in a file from its name.
  \item \textbf{Keep public surfaces small.} Mark only what other modules truly
        need with \code{pub}; keep helpers private. The less a module exposes, the
        freer you are to change its insides (Chapter~15).
  \item \textbf{Let data own its rules.} Put the operations that protect a struct's
        invariants in its methods, as the \code{Account} did in Chapter~9, so the
        rules live with the data they guard.
  \item \textbf{Depend on interfaces, not concrete types,} where you want
        flexibility (Chapter~11) --- it lets you swap implementations without
        rewriting the code that uses them.
\end{itemize}

\section{Coding style}

Style is not decoration; it is how you make code readable to the next person ---
who is often you, months later. Salam even ships a formatter, \code{salam fmt},
that rewrites your code into the canonical layout automatically, so you never argue
about spacing. Beyond formatting, a few habits carry most of the weight:

\begin{itemize}[leftmargin=1.4em]
  \item \textbf{Name for meaning.} \code{user\_count} and \code{is\_valid} tell the
        reader what a value \emph{is} and what a function \emph{answers}.
        Single-letter names earn their place only as short-lived loop indices.
  \item \textbf{Prefer immutability.} Reach for \code{mut} only when a value must
        change (Chapter~2). Every binding that cannot change is one fewer thing to
        track.
  \item \textbf{Keep functions small and single-purpose.} If you cannot name a
        function for the one thing it does, it is doing too many things
        (Chapter~6).
  \item \textbf{Comment the \emph{why}.} The code already shows \emph{what} it
        does; a comment is most valuable when it records the reasoning the code
        cannot.
  \item \textbf{Handle the unhappy path.} Check the \code{bool} a function returns,
        the \code{Option} that might be empty, the pointer that might be
        \code{null}. The cases you ignore are the bugs you ship.
\end{itemize}

\begin{tip}
Run \code{salam fmt} before you save or share code, and let it settle every
question of layout. Consistent formatting makes differences between versions show
real changes, not reshuffled whitespace, and frees your attention for what the code
actually does.
\end{tip}

\section{Common pitfalls, revisited}

A consolidated list of the traps this book has flagged --- worth a glance whenever
something behaves unexpectedly:

\begin{itemize}[leftmargin=1.4em]
  \item \textbf{Integer division truncates.} \code{7 / 2} is \code{3}; use a float
        operand for a fractional result (Chapter~3).
  \item \textbf{Literals are \code{i32}.} Storing one in a smaller or unsigned type
        needs an explicit \code{as} cast (Chapter~2).
  \item \textbf{Off-by-one with indices.} Arrays run \code{0} to \code{n-1}, while
        \code{repeat A to B} \emph{includes} \code{B} (Chapters~7 and~5).
  \item \textbf{Forgetting \code{free}.} Growable containers and manual allocations
        leak unless freed; pair each with a \code{defer} at creation
        (Chapter~18).
  \item \textbf{Unprotected shared state.} Two threads touching the same data
        without a lock race (Chapter~20).
  \item \textbf{Ignoring a result.} A \code{bool} or \code{Option} that signals
        failure is only useful if you check it (Chapter~14).
  \item \textbf{An \code{=} where you meant \code{==}.} Assignment versus
        comparison (Chapter~3).
\end{itemize}

\section{Where to go next}

The best way to consolidate everything in this book is to build something real ---
something a little beyond what you are sure you can do. A few projects that exercise
the whole language:

\begin{itemize}[leftmargin=1.4em]
  \item A \textbf{command-line tool} that reads a file, processes it, and writes a
        result --- exercising \code{io}, \code{str}, and the collections.
  \item A small \textbf{interpreter or calculator} that parses expressions ---
        exercising enums, structs, recursion, and error handling.
  \item A \textbf{simulation or game} (a grid world, a card game) --- exercising
        structs, methods, arrays, and \code{rand}.
  \item A \textbf{data tool} that reads JSON or CSV, tallies something with a
        \code{HashMap}, and reports --- exercising the standard library end to end.
\end{itemize}

Read the standard library's own source, too: it is written in Salam, and seeing how
\code{Vector} and \code{HashMap} are built is one of the best ways to deepen your
understanding. And when you hit something the language does not yet do, remember
the FFI --- the whole world of C is one \code{extern} block away.

\section{A closing word}

You began this book, perhaps, never having written a line of code; you end it
knowing a complete, modern programming language --- one you can write in your own
words. Programming is a craft, and like every craft it rewards practice more than
reading. So close the book and open an editor. Write a program that does something
you care about. Make it work, then make it clear, then make it fast --- in that
order. Argue with the compiler until you both agree. That is the whole of the
craft, and you now have every tool you need to practise it.

\bigskip
\noindent\emph{Salam, and happy coding.}

\section{Chapter summary}

\begin{itemize}[leftmargin=1.4em]
  \item Salam is fast by default; for more, choose good data structures, avoid
        needless allocation, and \emph{measure} before optimising.
  \item Structure a program by responsibility, keep public surfaces small, and let
        data own its rules.
  \item Good style --- meaningful names, immutability, small functions, ``why''
        comments, handled error paths --- is how code stays readable; \code{salam
        fmt} handles layout.
  \item Keep the list of common pitfalls in mind when something surprises you.
  \item Consolidate by building real projects and reading the standard library's
        source.
\end{itemize}

\section{Exercises}

\exercise Take a program you wrote for an earlier chapter and improve it: run
\code{salam fmt} on it, rename anything unclear, and add a comment explaining one
non-obvious decision.

\exercise Find a function in one of your programs longer than a screenful and split
it into smaller, well-named helpers.

\exercise Pick one project from ``Where to go next'' and sketch its modules: what
files, what \code{pub} functions, what data structures.

\exercise Review one of your programs against the ``common pitfalls'' list and fix
any you find.

\exercise Open a standard-library module (such as \code{collections}) and read one
of its types end to end. Write a paragraph on how it works.
